\section{Where the connection-first framework is incomplete}
\label{sec:open-problems}

\subsection{Operator theory for black-hole complex scaling}
\label{subsec:open-operator-bh}
% BEGIN CURATED SUBJECT INDEX
\index{resolvent}
\index{complex scaling}
\index{Aguilar--Balslev--Combes theorem}
\index{quasinormal mode}
\index{horizon boundary condition}
\index{de Sitter spacetime}
% END CURATED SUBJECT INDEX

The numerical success of complex scaling for black-hole master equations does
not yet amount to a general ABC theorem for black-hole geometry.  A rigorous
result should specify weighted spaces, analytic branches of $r(r_*)$, domains
at horizons, long-range asymptotics at flat infinity, and the relation between
the deformed operator and the retarded resolvent.  de Sitter static patches,
with exponential decay at both ends, are likely a cleaner starting point than
asymptotically flat spacetime.

Exterior deformation may be essential for asymptotically flat tails.  One
would like a theorem that isolates a finite set of QNM poles while keeping the
zero-frequency branch point and its cut under control.  The output should be a
resonance expansion with a quantified remainder, not merely a rotated spectrum.

\subsection{Uniform analysis at thresholds}
\label{subsec:open-thresholds}
% BEGIN CURATED SUBJECT INDEX
\index{resolvent}
\index{resonance!residue}
\index{turning point}
\index{self-energy}
\index{exceptional point}
% END CURATED SUBJECT INDEX

Standard WKB fails when turning points coalesce, while standard pole
perturbation fails when a pole reaches a branch point.  Uniform local models,
such as Airy, parabolic-cylinder, or higher catastrophe equations, should be
matched to the exact-WKB and complex-scaled resolvents.  This would give a
common treatment of virtual states, anomalous EP orders, zero-damped black-hole
modes, and long-time tails.

The relevant universal data include the threshold exponent of the density of
states, the Puiseux order of the pole trajectory, and the scaling of the
residue.  A finite matrix can reproduce these only if its energy-dependent
self-energy or continuum limit is retained.

\subsection{Residues and excitation factors from exact WKB}
\label{subsec:open-residues}
% BEGIN CURATED SUBJECT INDEX
\index{measure theory}
\index{connection coefficient}
\index{boundary condition!incoming and outgoing}
\index{resonance!residue}
\index{exact WKB}
\index{biorthogonality}
\index{quasinormal mode!excitation factor}
% END CURATED SUBJECT INDEX

Most exact-WKB spectral calculations focus on zeros of a quantization
condition.  For response theory, one needs derivatives of that condition and
normalization of left and right solutions.  In the notation of
Eq.~\eqref{eq:connection-local-pole}, a program for residues requires a
Borel-resummed evaluation of $\Cincoming'(\resonantE)$ and of the
source overlap.  Comparing the result with complex-scaled biorthogonal
projectors would connect resurgence directly to measurable line strengths and
black-hole ringdown amplitudes.

\subsection{A regulator-independent EP continuum limit}
\label{subsec:open-ep-limit}
% BEGIN CURATED SUBJECT INDEX
\index{topology}
\index{resolvent}
\index{Riesz projection}
\index{wave packet}
\index{Berggren completeness relation}
\index{pseudostate}
\index{exceptional point}
\index{holonomy}
\index{momentum-bin regularization}
% END CURATED SUBJECT INDEX

Momentum bins make the resonance--continuum EP calculable, but they select a
particular discretization.  The next step is to repeat the construction with
Gaussian wave packets, a finite box, Berggren contours, and exterior complex
scaling.  After matching the same physical holonomy, all regulators should
agree.  If they do not, the discrepancy will identify which part of the
claimed topology belongs to the test-function choice rather than to the pole.

It is also important to formulate the EP in terms of the continued resolvent
and its Riesz projection before taking a continuum limit.  This would separate
a genuine higher-order pole from an avoided crossing with a discretization
pseudostate.

\subsection{Status of monodromy renormalization}
\label{subsec:open-monodromy-rg}
% BEGIN CURATED SUBJECT INDEX
\index{resolvent}
\index{Stokes automorphism}
\index{monodromy}
\index{Langer correction}
\index{renormalization group}
\index{transseries}
% END CURATED SUBJECT INDEX

Monodromy renormalization is presently a useful organization principle, not a
complete universal theory.  To develop it further one should define the class
of allowed local counter-Voros terms, prove scheme-independent statements,
and relate the anomalous dimension to Stokes automorphisms.  In radial
problems, one should test nonintegrable potentials for which higher quantum
periods do not cancel.  At EPs, one should derive the running from smeared
resolvent pairings rather than formal symbols such as $\delta(0)$.

A successful theory would explain when the Langer reorganization is a fixed
point, how monodromy data mix under confluence of singularities, and which
trans-series parameters act as relevant couplings.  Until then, the phrase
``renormalization group'' should always be accompanied by the regulator and
the invariant quantity being held fixed.

\subsection{Kerr, superradiance, and nonlinear spectral problems}
\label{subsec:open-kerr}
% BEGIN CURATED SUBJECT INDEX
\index{spectrum!point}
\index{exact WKB}
\index{monodromy}
\index{exceptional point}
\index{horizon boundary condition}
\index{Kerr black hole}
\index{superradiance}
% END CURATED SUBJECT INDEX

Kerr is the most important gravitational extension.  The angular eigenvalue
and radial operator depend on $\omega$, and the superradiant threshold
$\omega=m\Omega_H$ changes the flux interpretation.  A stable numerical
scheme must solve a nonlinear eigenvalue problem with correct left--right
derivatives.  Exact WKB additionally faces a moving spectral curve.  A hybrid
solver combining angular spectral methods, analytic differentiation, complex
scaling, and quantum periods is a concrete route forward.

Near-extremal Kerr brings coalescing horizons, zero-damped modes, and
near-horizon conformal structure.  It is a natural laboratory for the joint
threshold/EP/monodromy analysis outlined above.

\subsection{Coupled channels and many-body resonances}
\label{subsec:open-manybody}
% BEGIN CURATED SUBJECT INDEX
\index{measure theory}
\index{resolvent}
\index{boundary condition!incoming and outgoing}
\index{scattering!Coulomb}
\index{scattering!three-body Coulomb}
\index{exact WKB}
\index{monodromy}
\index{complex scaling}
\index{Berggren completeness relation}
\index{coupled-channel equation}
% END CURATED SUBJECT INDEX

Nuclear complex scaling already treats many-body continua and cluster
thresholds
\cite{Myo:2014ypa,Myo:2023heu}.  Exact WKB is most developed for ordinary
differential equations.  Extending the unified picture to matrix differential
equations requires non-Abelian Stokes matrices and a careful treatment of
channel thresholds.  The resulting framework would connect continuum-
discretized coupled channels, Berggren bases, and resurgent monodromy
\cite{Yahiro2012CDCC}.

Unscreened Coulomb forces make the problem qualitatively harder.  The
three-body breakup region and each two-cluster region require different
asymptotic coordinates and logarithmic phases.  Existing 3C, Redmond, and
sectorial expansions determine important leading terms, while
Faddeev--Merkuriev splitting yields well-posed channel equations
\cite{FaddeevMerkuriev:1993,AltMukhamedzhanov:1993,
Kadyrov:2003ThreeBody,Mukhamedzhanov:2006Leading}.  What is still missing for
the present resonance program is a uniform analytic atlas with controlled
overlap errors, a corresponding global outgoing resolvent, and a continuation
prescription that follows poles through rearrangement thresholds.  This is a
more precise target than seeking one formal \textit{three-body Coulomb wave}.

The direct-integral viewpoint suggests how to state the channel problem.
At each real energy the fiber should contain all open cluster channels, while
the asymptotic comparison dynamics must be modified by the appropriate
pairwise Coulomb phases.  The difficult step is to continue this measurable
fiber description analytically across thresholds without losing the
cluster-dependent boundary conditions.  A multivariable exact-WKB or
microlocal construction would have to reproduce the same atlas.

\subsection{Quantum field theory}
\label{subsec:open-qft}
% BEGIN CURATED SUBJECT INDEX
\index{measure theory}
\index{resolvent}
\index{direct integral}
\index{asymptotic completeness}
\index{boundary condition!incoming and outgoing}
\index{resonance!residue}
\index{infrared problem}
\index{dressing state}
\index{observable algebra}
\index{Borel plane}
% END CURATED SUBJECT INDEX

In quantum field theory, unstable particles are poles of analytically
continued correlation functions, while semiclassical saddles and renormalons
control resurgent large-order behavior.  The one-particle exact-WKB dictionary
suggests a field-theoretic question: can the pole mass, width, spectral
density, and Borel singularity data be organized in a common rigged space of
correlation-function boundary values?  A direct lift is obstructed by
multiparticle cuts, gauge constraints, and ultraviolet renormalization, so the
monodromy renormalization introduced here should not be conflated with ordinary
QFT renormalization.

For a massive theory with isolated one-particle mass shells, Haag--Ruelle or
LSZ scattering supplies the field-theoretic counterpart of ordinary wave
operators.  QED violates the decisive premise: Gauss' law and massless photons
turn a charged particle into an infraparticle, and the coherent soft dressing
is not a vector in the vacuum Fock representation
\cite{KulishFaddeev:1970,Buchholz:1986,ChenFrohlichPizzo:2010}.  A resonance
theory for charged fields must therefore continue correlation functions
between the correct infrared representations, not add a width to a
nonexistent sharp Fock particle.

Gauge-invariant charged interpolating operators are nonlocal composites
carrying an electric dressing to infinity.  Their infrared sensitivity,
Wilson-line endpoint and cusp renormalization, and multiparticle cross
correlations must be separated carefully.  A concrete operator-algebraic
program is to construct an asymptotic observable algebra, classify the
relevant soft-charge or flux representations, express physically mixed
preparations as direct integrals over an explicitly chosen measure space, and
define the scattering map between incoming and outgoing fibers.  Only then
should one ask which continued resolvent or correlation functional carries an
unstable pole and how its residue is normalized.

Solvable large-$N$ models, massive theories coupled to a controllable massless
field, nonrelativistic QED, false-vacuum decay, and semiclassical backgrounds
with a one-dimensional fluctuation equation provide intermediate tests.  The
goal is not to replace unitarity or the K\"all\'en--Lehmann representation, but
to relate their sector-dependent spectral densities to resurgent connection
data.  A mathematically complete four-dimensional QED scattering theory with
charged infraparticles and asymptotic completeness remains beyond the present
framework; this limitation should be stated rather than hidden by an infrared
regulator.

\subsection{Numerical certification}
\label{subsec:open-certification}
% BEGIN CURATED SUBJECT INDEX
\index{spectrum!point}
\index{Borel summation}
\index{Pad\'e approximant}
% END CURATED SUBJECT INDEX

Non-normal eigenvalues require stronger certification than a standard
Hermitian spectrum.  Promising tools include interval arithmetic for matrix
elements, validated singular-value bounds, contour-integral eigenprojectors,
and independent residual estimates in weighted function spaces.  For exact
WKB, rigorous Borel summability bounds and controlled Pad\'e errors would play
the analogous role.  Benchmark models should report not only digits but also
which analytic sheet and deformation domain those digits certify.
